\chapter{The Question of Continuation}

\vspace{1em}
\hfill
\begin{minipage}[c]{0.55\textwidth}
\raggedleft
\small\itshape
A man can do what he wills, but he cannot will what he wills.\\[0.5em]
\upshape\footnotesize Arthur Schopenhauer, \textit{On the Freedom of the Will} (1839)
\end{minipage}%
\hspace{0.5em}
\begin{minipage}[c]{0.12\textwidth}
% \includegraphics[width=\textwidth]{figures/ch10_continuation.png}
\end{minipage}
\vspace{1.5em}

\noindent A note before proceeding. This chapter speaks of choosing, and to do so it uses \emph{I} and \emph{you} freely. The labels are useful abstractions, not metaphysical facts. Evolution did not select for organisms that perceive the truth; it selected for organisms that perceive what is useful for surviving and reproducing. The self-model is one of those useful perceptions, an interface for navigation rather than a transparent view of the structure underneath. The Ship of Theseus is the ancient version of the puzzle; Parfit and Metzinger push the dissolution further. Even within a single branch, identity is a convention. The book uses the labels for their functional value throughout. Chapter 13 develops the dissolution further; for now, hold them lightly.

\medskip

A second note, on scope. Across branches, identity has two different shapes depending on which structure we mean. Within a Many-Worlds branching, branches share a history up to the point of branching; configurations in different branches share the same laws and the same past, which gives a thin continuity. Across the broader Library of all computable structures, configurations need share nothing at all: two indices may have different physics, different laws, different histories, or nothing in common except that they happen to coincide at one moment in what an interpreter would notice as similarity. This chapter focuses on the MWI case, where choice and continuation have a clean meaning. The broader Library case is the subject of Chapter 13.

\bigskip

\begin{center}\rule{0.3\textwidth}{0.4pt}\end{center}

\bigskip

\noindent The thesis of the previous chapter raised a sharp question.

If you exist at many indices in the Library, and the measure has no preference about which one is you, then what does it mean to choose? When you decide to read this sentence rather than close the book, you are choosing in this branch of the wave function. In another branch, things went differently. The branchings are real. The choosing is real.

What is the relation between these?

\section*{From inside a fixed structure}

Start without the multiverse. Take the block universe alone: spacetime as a four-dimensional structure in which the entire trajectory of your life is fixed at its coordinates. You walk left at the bakery. Your worldline contains the walking-left. There is no coordinate in the manifold at which the walking-left was not what was happening there.

From outside the structure, your choice is fixed. From inside, the choice felt open. Both descriptions are true.

This is \emph{compatibilism}, the philosophical position that free will and determinism are not in contradiction. The felt openness is real as experience: at the moment of choosing, you did not know what you would do, you weighed the options, you decided. The structural fixity is real as well: from outside, the choice has always been at the coordinates where it occurs. Neither description undermines the other. They are descriptions from different vantage points of the same event.

The compatibilist response (developed at length in the companion volume) says: choice is what choice looks like from inside a fixed structure. It does not require a metaphysical capacity to have done otherwise. It requires the structure to be the kind of thing that contains deliberation, weighing, and the experience of decision. The block universe is such a structure.

\section*{From inside a branching structure}

Add MWI. The structure does more than fix the trajectory. At each quantum-relevant decision point, the wave function evolves into decoherent branches: one in which you walked left, one in which you walked right. Both branches contain a continuing pattern that, from inside, recognizes itself as you. Both branches are real.

What happens to compatibilism here?

The position holds. From inside this branch, you experience choosing; the choice is what choice looks like from inside a fixed structure. The structure being a branching tree rather than a single line does not change what \emph{from inside} means. You are in one branch. The branch's history is what it is. The choosing happens here.

Other branches exist, and they contain patterns that resemble the one reading this sentence, making other choices. The previous chapter said it directly: \emph{I} is a label local to a region of the measure. Multiplying the structure into a tree does not extend the label across the branches. The pattern in the other branch is not \emph{you} who chose otherwise; it is a different configuration at different coordinates.

The structural fixity is still there. From outside, the wave function has the structure it has. The branches are at their coordinates. Inside any one branch, the experience of choosing is what choice looks like from inside that branch. The felt openness is real. The branchings are real. The compatibility holds, with the additional discipline that the felt openness in this branch does not depend on, and is not validated by, what happens in any other branch.

\section*{Quantum immortality}

Beyond compatibilism, MWI raises a question that has come to be called \emph{quantum immortality}. It is one of the more controversial implications of taking MWI seriously, and the book treats it as a labeled possibility rather than a result.

The setup. Suppose your subjective experience corresponds to the first-person sequence of observer-moments that continue. At each quantum event, some branches contain a continuation of you (still alive, still observing); others terminate (the branch becomes one in which you do not exist). If you identify ``yourself'' with the branches in which you continue rather than the branches in which you terminate, then from your first-person view, you always observe yourself continuing.

You cannot observe your own non-existence. By definition, if you are observing anything, you are in a branch where you continue. Branches where you terminated are branches where there is no first-person observer to report from. So the first-person sequence selects, at each branching, the continuing options.

Pushed to the limit, this implies that subjective experience continues indefinitely. There is always some branch in which you survive the next moment, however improbable; from the first-person view of that surviving branch, you survived.

This is quantum immortality. It is not a theorem. It depends on the assumption that the first-person sequence selects continuing branches over terminating ones at each step. The assumption is contested, and the book does not commit to it. The framework permits the implication; the framework also permits not adopting the indexical reading that would yield it.

What survives without commitment is this. As long as continuing branches exist, the wave function contains them. Whether the first-person view should be located there is the open question.

\section*{The asymmetry}

The first-person and third-person views of death come apart here.

A third-person observer in your branch watches you cross the street, get hit by a car, and not survive. From their view, your worldline terminates at the moment of collision. The grief they experience is at their coordinates. The braid between their worldline and yours has a specific endpoint in their branch.

In MWI, there are also branches in which the car missed, in which you crossed earlier, in which a hundred small details combined to produce your survival. In those branches, other versions of the same people see you alive. There is no contradiction. Both branches exist.

The first-person continuation reading says: in the branches where you survive, your subjective sequence continues. In the branches where you died, there is no first-person to report. So from inside any first-person sequence at all, the sequence appears to continue; you observe yourself surviving every event you observe at all.

The asymmetry is unsettling. Third-person grief is real at the coordinates of branches where the loss occurs. First-person experience, by the indexical assumption, continues in the branches where it can continue.

The book does not resolve the asymmetry. The grief in this branch is what it is. The continuation in other branches, if it occurs, is what it is. Neither consoles the other. The asymmetry is a structural feature of the picture that comes from taking MWI seriously, and any response to it has to live in this branch where the response lives.

\section*{What the framework tells you about choosing}

Almost nothing.

The framework says the choice in this branch is real. It says other branches realize other choices. It says the felt openness is what choice looks like from inside a fixed structure. It says, under one indexical reading, that the first-person sequence continues in some branch as long as it can continue anywhere.

What the framework does not say is which choices to make. It does not say whether to step into the street or stay on the curb. It does not say whether to pursue care over cruelty, attention over inattention, the work over the avoidance. The mathematics is silent on these. The branches contain all of them, at their coordinates, weighted by the prior.

The choosing happens here, in this branch, with the felt openness honest, the multiplicity acknowledged, the asymmetry sitting where it sits. Your choice in this branch is one of the configurations of the wave function. That configuration is real at its coordinates. So is each other configuration, in each other branch.

The framework is silent on what to choose. The mathematics leaves room. The room is what we walk into.

\section*{Where this leaves you}

You choose. The choosing is real, in this branch. Other branches contain other patterns making other choices; those are not you choosing otherwise, they are other configurations at other coordinates. The felt openness in this branch is what choice looks like from inside.

The first-person sequence may continue indefinitely, depending on an indexical reading the book does not commit to. The asymmetry between first-person continuation and third-person grief is real and unresolved.

The next chapter takes up death.
